\section*{Chapitre \ref{ch:hf:automated-options-market-making} --- Market making automatisé d'options}

\begin{solution}{exo:hf:automated-options-market-making:1}
Un demi-point de chaque côté~: $2\times0.5\times0.114=\$0.114$, soit environ 11 cents.
\end{solution}

\begin{solution}{exo:hf:automated-options-market-making:2}
L'option d'achat vaut désormais environ 21.00~\$~: $0.99\times10\times100=\$990$.
\end{solution}

\begin{solution}{exo:hf:automated-options-market-making:3}
Sinon, chaque ordre légitime comportant autant de jambes déclenche la protection et retire toutes les cotations, ce qui fait perdre les exécutions qui arrivent pendant que celles-ci sont retirées.
\end{solution}

\begin{solution}{exo:hf:automated-options-market-making:4}
Les premières cibles sont des options profondément dans la monnaie qui gagnent toutes à peu près la même chose (leur delta vaut un)~: chaque exécution ajoute environ 1\,000~\$.
\end{solution}

\begin{solution}{exo:hf:automated-options-market-making:5}
Le balayage continue de toucher les autres cotations de l'échéance jusqu'à ce que l'annulation prenne effet~: à raison d'une série toutes les 5 microsecondes, 50 microsecondes représentent dix
exécutions supplémentaires.
\end{solution}

\begin{solution}{exo:hf:automated-options-market-making:6}
Les preneurs de liquidité qui choisissent le moment de leurs exécutions négocient lorsque la cotation du teneur de marché lui est le moins favorable~; les largeurs doivent tarifer la prévisibilité à court
terme de la fair value, ou les cotations doivent se déplacer plus vite que le timing des preneurs.
\end{solution}

\begin{solution}{exo:hf:automated-options-market-making:7}
25\,900~\$ avec un seuil de 13~; 164\,300~\$ sans protection (151 séries mal valorisées).
\end{solution}

\begin{solution}{exo:hf:automated-options-market-making:8}
Les balayages qu'il a observés étaient petits~; un mouvement de 1~\% avec un seuil de quarante coûterait 38\,700~\$ dans le modèle. Fixer le seuil d'après le pire
balayage plausible, et non d'après l'historique.
\end{solution}

\begin{solution}{pb:hf:automated-options-market-making:1}
\begin{enumerate}
\item Voir la \cref{def:hf:automated-options-market-making:live}~; le sous-jacent fait évoluer la surface en continu, tandis que le smile est réajusté à une cadence plus lente.
\item À 0.025 point de volatilité près, avec 0.2 point de bruit par cotation.
\item Voir la \cref{def:hf:automated-options-market-making:mass}.
\item 11 cents pour l'option d'achat à la monnaie à 30 jours~; d'un cent à 20 cents.
\item Paramètres de notionnel, de volume, de nombre et de pourcentage de cotation sur des fenêtres allant de 100 ms à une journée~; le trading s'arrête dès que l'un d'eux est atteint.
\item Un mouvement de 1~\%, 150 cotations périmées, une touche toutes les 5 microsecondes~; ordres ordinaires toutes les 5 secondes, de 1 à 12 jambes, 4~\$ par exécution.
\item 4\,950~\$ et 1\,052~; 9\,899~\$ et 326~; 12\,867~\$ et 3~; 19\,773~\$ et 0.
\item Le balayeur prend en premier les plus gros écarts de valorisation, car la protection peut l'arrêter à tout moment.
\item Position courte d'environ 13\,000 actions avec un seuil de 13~; la couverture empêche les mouvements suivants d'ajouter à la perte.
\item Le risque d'inventaire a un effet de premier ordre sur les prix des options, plus grand que l'information asymétrique dans l'impact sur les prix.
\item Voir la \cref{def:hf:automated-options-market-making:strike}.
\item 16\,500~\$ sans protection~; 9\,900~\$ avec une limite de 1\,000~\$ à 50 microsecondes, 6\,700~\$ à 5.
\item De 1\,980~\$ pour deux exécutions à 79\,370~\$ sans protection, soit environ 1\,000~\$ par exécution~; les seuils inférieurs à douze perdent 326 à 2\,441 exécutions par jour,
treize n'en perd que trois.
\item Juste au-dessus du plus grand ordre légitime à plusieurs jambes~: treize ici.
\item Un balayage sur plusieurs bourses à la fois rencontre plusieurs compteurs, chacun ne comptant que ses propres exécutions.
\item La bourse compte les exécutions, et non le risque~: le vega et le gamma s'accumulent d'une manière qu'un comptage ne peut pas voir.
\item Les exécutions dont le timing est choisi rapportent moins de 40~\% des fourchettes conventionnelles.
\item Market making par \omterm{def:hf:automated-options-market-making:mass}{cotation de masse}.
\item À partir de la perte que la firme accepte sur un saut de volatilité de taille plausible, et de la vitesse de mise en œuvre de la protection.
\item Pertes dues aux balayages des traders informés, face aux exécutions perdues à cause des ordres ordinaires.
\end{enumerate}
\end{solution}

\begin{solution}{iq:hf:automated-options-market-making:1}
La surface évolue avec le sous-jacent~; chaque cotation est recalculée et renvoyée en \omterm{def:hf:automated-options-market-making:mass}{cotations de masse}~; tant qu'elles ne sont pas parvenues, les cotations périmées peuvent être balayées et la
protection peut les retirer~; le delta issu des exécutions éventuelles est couvert~; le smile est réajusté si le mouvement l'a modifié.
\iqlookfor{l'enchaînement et la protection.}
\end{solution}

\begin{solution}{iq:hf:automated-options-market-making:2}
Décodage des données de marché, mise à jour de la surface, calcul de 5\,000 valeurs théoriques et des largeurs (vectorisé), encodage des messages par blocs selon les limites de la bourse, envoi~; l'essentiel du temps
part au calcul et à l'encodage, donc il faut précalculer les sensibilités et mettre à jour linéairement.
\iqlookfor{vectorisation et mises à jour linéaires.}
\end{solution}

\begin{solution}{iq:hf:automated-options-market-making:3}
Faire évoluer le smile avec le sous-jacent selon la règle, et le réajuster lorsque les cotations ou les transactions du marché s'écartent du smile déplacé au-delà de leur bruit, ou
après un nombre fixé de mises à jour.
\iqlookfor{résidus par rapport à un seuil de bruit.}
\end{solution}

\begin{solution}{iq:hf:automated-options-market-making:4}
Si les déclenchements étaient des balayages ou des ordres légitimes à plusieurs jambes, les pertes dans chaque cas, le seuil face aux plus gros ordres, et le timing par rapport
aux mouvements du sous-jacent.
\iqlookfor{séparer les déclenchements informés des légitimes.}
\end{solution}

\begin{solution}{iq:hf:automated-options-market-making:5}
Les acheteurs de puts sont plus souvent informés ou couvrent un risque de krach, le vega de l'aile des puts est plus difficile à couvrir, et l'inventaire y porte un risque de queue de distribution.
\iqlookfor{risque et flux asymétriques.}
\end{solution}

\begin{solution}{iq:hf:automated-options-market-making:6}
La perte vaut $\sum_{i=1}^{\min(n,N)}e_i$ multipliée par la taille et le multiplicateur pour $N$ cibles~; avec des edges égaux elle est linéaire en $n$, avec des edges décroissants
concave.
\iqlookfor{sommes partielles des edges ordonnés.}
\end{solution}
